%!TEX program = pdflatex
%
% ── Journal of Chemical Information and Modeling (ACS) submission style ────
\documentclass[journal=jcisd8]{achemso}

% ── Core packages ──────────────────────────────────────────────────────────
\usepackage[utf8]{inputenc}
\usepackage[T1]{fontenc}
% Times-compatible fonts matching ACS/JCIM published style (12 pt base via achemso)
\usepackage{newtxtext}   % Times Roman body + sans for headings
\usepackage{newtxmath}   % Times-compatible math (includes AMS symbols; replaces amssymb)
\usepackage{amsmath,mathtools}
\usepackage{booktabs,multirow,makecell,array,tabularx}
\usepackage{threeparttable}
\usepackage{graphicx}
\graphicspath{{figure/}}
\usepackage{xcolor}
\usepackage{float}
\usepackage[section]{placeins}
\usepackage[colorlinks=true, linkcolor=blue!60!black,
  citecolor=green!50!black, urlcolor=blue!70!black]{hyperref}
\usepackage{cleveref}
\usepackage{enumitem}
\usepackage{subcaption}
\usepackage{siunitx}
\usepackage{microtype}
\raggedbottom
\setlength{\parindent}{1em}   % JCIM: first-line indent on body paragraphs

% \doi not provided by achemso when using thebibliography (only via .bbl)
\providecommand*{\doi}[1]{\href{https://doi.org/#1}{\texttt{doi:#1}}}

\SectionNumbersOn
\setcounter{secnumdepth}{3}

% Affiliation markers: arabic numerals (1, 2 …); corresponding-author marker: *
% Arrangement: affil number first, then * → e.g. "1,*" not "*,1"
\makeatletter
% (1) Symbol: plain * for slot 0 (email/corresponding), arabic numeral otherwise
\def\acs@author@fnsymbol#1{\ifnum#1=\z@ *\else\number#1\fi}
% (2) Build author footnote string with affil BEFORE star
\def\acs@author@list@main{%
  \advance\@tempcnta\@ne\relax
  \def\acs@author@footnotes{}%
  \acs@author@list@and
  \space
  \@nameuse{@author@\@roman\@tempcnta}%
  \acs@author@list@comma
  \acs@author@affil
  \acs@author@affil@also
  \acs@author@affil@alt
  \acs@author@star
  \ifx\@empty\acs@author@footnotes\else
    \textsuperscript{\acs@author@footnotes}%
  \fi
  \ifnum\@tempcnta<\acs@author@cnt\relax
    \expandafter\acs@author@list@main
  \fi
}
% (3) Star@aux: append ,* after the affil number already in \acs@author@footnotes
\def\acs@author@star@aux{%
  \protected@edef\acs@author@footnotes{%
    \acs@author@footnotes
    \ifnum\acs@affil@cnt>\@ne\relax
      ,%
    \else
      \ifnum\acs@affil@alt@cnt>\z@\relax
        ,%
      \fi
    \fi
    \acs@author@fnsymbol{\z@}%
  }%
}
\makeatother

\DeclareMathOperator{\softmax}{softmax}
\DeclareMathOperator{\softplus}{softplus}
\DeclareMathOperator{\ReLU}{ReLU}
\newcommand{\vect}[1]{\boldsymbol{#1}}
\newcommand{\mat}[1]{\mathbf{#1}}

% ── Author / affiliation / title (ACS achemso style) ─────────────────────
\title{MOSAIC-GNN: A Multi-Graph Neural Network for High-Accuracy
  Molecular Formula Assignment from MS1-Only Mass Spectrometry Data}

\author{Supaporn Klabklaydee}
\affiliation{Department of Civil and Environmental Engineering,
  School of Environment and Society,
  Institute of Science Tokyo, Tokyo, Japan}
\email{klabklaydee.s.aa@m.titech.ac.jp}

\author{Nopphakorn Subsa-Ard}
\affiliation{Department of Computer Engineering,
  Faculty of Engineering,
  King Mongkut's University of Technology Thonburi,
  Bangkok 10140, Thailand}

\author{Suppanut Chankhow}
\affiliation{Department of Civil and Environmental Engineering,
  School of Environment and Society,
  Institute of Science Tokyo, Tokyo, Japan}

\author{Nattapak Darumas}
\affiliation{Department of Eye and Vision Sciences,
  Institute of Life Course and Medical Sciences,
  University of Liverpool, Liverpool, United Kingdom}

\author{Gu Yuheng}
\affiliation{Department of Civil and Environmental Engineering,
  School of Environment and Society,
  Institute of Science Tokyo, Tokyo, Japan}

\author{Manabu Fujii}
\affiliation{Department of Civil and Environmental Engineering,
  School of Environment and Society,
  Institute of Science Tokyo, Tokyo, Japan}
\email{fujii.m.ah@m.titech.ac.jp}

\keywords{molecular formula assignment, graph neural network, mass spectrometry,
  metabolomics, dissolved organic matter, energy-based ranking, cheminformatics,
  uncertainty quantification}

% ═══════════════════════════════════════════════════════
\begin{document}

\begin{abstract}
Accurate molecular formula assignment from high-resolution mass spectrometry (MS) data
is a critical computational bottleneck in untargeted metabolomics and
environmental chemistry.
Existing methods either require MS2 fragmentation spectra or achieve limited accuracy
when restricted to MS1 data alone.
We present MOSAIC-GNN
(\textbf{M}ixture-\textbf{O}ptimised \textbf{S}pectral--\textbf{A}dduct--\textbf{I}sotope
\textbf{C}o-Inference Graph Neural Network), a multi-graph neural network that assigns
molecular formulas from MS1 data by co-inferring over four complementary graph
representations: a Spectral Mass-Difference Graph (SMDG), an Isotopologue Pattern
Graph (ISO), a Formula Candidate Lattice (FCL), and an Adduct Consistency Graph (ADC).
Candidates are ranked via per-feature energies (log-softmax over negated scores),
using energy-based \emph{ranking} rather than a generative energy-based model.
Evaluated on 4,047 dissolved organic matter (DOM) spectral features from three IHSS
aquatic fulvic acid standards (Pahokee River FA, Suwannee River FA\,2, Suwannee River
FA\,3) that are compositionally similar to the training data, the best-performing
configuration (M10: FCL-14 + Bethe--Weizs\"{a}cker physics features) achieves
\textbf{97.68\% Top-1 accuracy} (MRR\,=\,0.983), surpassing the best MS1-only
baseline (BUDDY, 63.4\%) by 34.3\,pp and exceeding SIRIUS\,5 with MS2 fragmentation
data (87.3\%) by 10.4\,pp.
As a more demanding cross-domain test, \textbf{86.40\% Top-1 accuracy} is achieved on
3,235 structurally diverse GNPS library entries spanning pharmaceutical, dietary, and
natural product compounds \emph{without any fine-tuning}---establishing the primary
cross-domain generalisability metric.
PFAS transfer learning achieves 78.9\% Top-1 (Top-5: 92.0\%) after partial-weight
fine-tuning with a formula-bypass pathway.
Ablation analysis identifies spectral neighbourhood context (SMDG, $+$5.39\,pp) and
FiLM instrument-type conditioning ($+$4.20\,pp) as the primary accuracy drivers.
Element-class stratification reveals near-perfect performance for CHO (99.8\%) and
CHON (94.2\%) compounds, while CHONS (18.9\%) and CHOPN (2.2\%) expose a fundamental
candidate-space bottleneck for heteroatom-rich formulas.
MOSAIC-GNN demonstrates that physics-informed multi-graph co-inference achieves
state-of-the-art MS1-only formula assignment.  The complete implementation
(2,056,119 parameters for M10) is publicly available together with a
web platform (Next.js\/FastAPI, supporting drag-and-drop mzML\/CSV upload)
and a standalone cross-platform desktop application,
enabling non-expert users to run formula assignment without programming.
\end{abstract}

\maketitle

\noindent\textbf{Keywords:} molecular formula assignment, graph neural network, mass spectrometry,
metabolomics, dissolved organic matter, energy-based ranking, cheminformatics, uncertainty quantification

\section{Introduction}
\label{sec:intro}

\begin{figure}[tp]
  \centering
  \includegraphics[width=\linewidth]{overview.png}
  \caption{System overview of MOSAIC-GNN. From a raw high-resolution MS1 peak list,
    four complementary graphs are constructed: (A) Spectral Mass-Difference Graph (SMDG)
    encoding mass-difference relationships between features; (B) Isotopologue Graph (ISO)
    capturing $M$/$M{+}1$/$M{+}2$ isotope patterns; (C) Formula Candidate Lattice (FCL)
    enumerating chemically valid candidates per feature; and (D) Adduct Consistency
    Graph (ADC) linking features related by common adducts.  The four-graph co-inference
    model assigns the top-ranked molecular formula with confidence score and elemental
    class label as output.}
  \label{fig:overview}
\end{figure}

High-resolution mass spectrometry (HRMS) has become the indispensable workhorse of untargeted
metabolomics, exposome research, and environmental chemistry.  A modern Fourier Transform Ion
Cyclotron Resonance (FT-ICR) or Orbitrap instrument routinely measures thousands of spectral
features per sample at sub-ppm mass accuracy, yet the fundamental problem of confidently
inferring the correct molecular formula for each observed \(m/z\) remains deceptively
difficult~\cite{Kind2007,Hufsky2020}.

\FloatBarrier
\subsection{The Molecular Formula Assignment Problem}

Given a set of observed $m/z$ values, the goal is to assign each peak its correct molecular formula
$\mathrm{C}_c\mathrm{H}_h\mathrm{O}_o\mathrm{N}_n\mathrm{S}_s\mathrm{P}_p\cdots$
that is consistent with the observed mass.  Even with 1~ppm mass accuracy, the number of
chemically plausible candidates grows rapidly with mass: at 400~Da, several dozen CHO-class
formulas may match within tolerance, and each additional heteroatom (N, S, P) exponentially expands
the combinatorial search space~\cite{Kind2007}.  The problem is further compounded by
adduct formation, in-source fragmentation, and the isotopologue envelope that partially
overlaps the monoisotopic peak.

Traditional approaches relied on rule-based enumeration with physical plausibility filters.
The ``Seven Golden Rules'' of Kind and Fiehn~\cite{Kind2007} provide a set of element-ratio
constraints (H/C, O/C, N/C, S/C, P/C) and nitrogen-parity checks that substantially prune
the candidate set.  Tools such as MFAssignR~\cite{Tolić2017} implement these rules with
Kendrick mass defect filtering for homologous series assignment in DOM and implement automated
formula assignment pipelines for FT-ICR MS data.  While effective for pure CHO compound
classes, these rule-based methods struggle when heteroatom-rich formulas (CHONS, CHOPN)
produce large candidate sets and when compound classes outside established DOM heuristics
are encountered.

\subsection{Fragmentation-Dependent Methods and Their Limitations}

SIRIUS~\cite{Duehrkop2019,Duehrkop2015} represents the current state of the art for formula
assignment in untargeted metabolomics, combining isotopologue scoring, fragmentation tree
construction, and machine learning (CANOPUS~\cite{Dührkop2021}) to assign formulas with high
confidence.  However, SIRIUS fundamentally requires tandem mass spectrometry (MS/MS) data:
its fragmentation tree algorithm cannot operate without at least a few MS2 fragment ions.
In practice, a majority of features detected in complex mixtures lack MS2 spectra because
data-independent acquisition does not achieve comprehensive fragmentation coverage, and
data-dependent acquisition biases toward abundant peptides at the expense of low-abundance
metabolites~\cite{Aron2020}.

BUDDY~\cite{Guo2023} partially addresses this gap via bottom-up metabolite annotation that
assembles formula candidates from neutral loss patterns but still achieves substantially lower
Top-1 accuracy (63.4\% in our benchmark, Section~\ref{sec:benchmark}) compared to
fragmentation-tree methods.  Neither BUDDY nor SIRIUS was designed with DOM or complex
environmental matrices in mind, where thousands of structurally similar CHO/CHON features
at similar masses create an especially challenging disambiguation task.

\subsection{Machine Learning for Molecular Formula Assignment}

The rise of graph neural networks (GNNs) in cheminformatics~\cite{Gilmer2017,Reif2019}
has opened new avenues for formula inference.  Several studies have exploited GNNs for
molecular property prediction~\cite{Yang2019}, retention time prediction~\cite{Bouwmeester2019},
and collision cross-section prediction~\cite{Mesleh1996}, but GNN-based formula
\emph{assignment} from raw MS1 spectra has remained largely unexplored.

Key challenges specific to this task include: (i) the molecular formula is not directly
encoded as a graph---the input is a raw spectrum---requiring purpose-built graph
representations; (ii) the number of candidate formulas per feature varies enormously
(1 to 500+), demanding a flexible ranking architecture; (iii) multiple complementary
information channels (isotopologue patterns, mass differences, adduct relations) must be
jointly exploited; and (iv) domain adaptation across instrument types and chemical classes
is essential for practical utility.

\subsection{Contributions}

In this work, we address these open challenges with MOSAIC-GNN.  Our main contributions are:

\begin{enumerate}[leftmargin=1.5em, label=\arabic*.]
  \item \textbf{Multi-graph MS1 representation.}
    We design four purpose-built graph representations that jointly encode spectral context
    (SMDG), isotopologue physics (ISO), formula-level relationships (FCL), and adduct
    consistency (ADC), enabling rich information fusion from MS1 data alone.

  \item \textbf{Physics-informed neural architecture.}
    We embed domain knowledge at multiple levels: the Seven Golden Rules act as hard
    constraints; differentiable Roepstorff isotope ratio prediction provides physics-informed
    supervision; Bethe--Weizs{\"a}cker nuclear mass-defect features enrich formula representations;
    and a learnable adduct consistency penalty enforces neutral-mass coherence.

  \item \textbf{Energy-based ranking with GIN, bypass, and interaction pathways.}
    The Formula GIN Scorer uses three complementary pathways (GIN message passing,
    formula bypass, and feature--formula dot-product interaction) to score candidate
    formulas as energies---using energy-based \emph{ranking} (log-softmax over
    negated energies per feature) rather than a generative energy-based model in
    the Lecun--Hinton sense---with per-feature InfoNCE contrastive training.

  \item \textbf{Evidential uncertainty quantification.}
    Normal-Inverse-Gamma posteriors provide calibrated aleatoric and epistemic uncertainty
    per candidate, enabling quality-control filtering in downstream analyses.

  \item \textbf{State-of-the-art performance and generalisability.}
    MOSAIC-GNN (M10) achieves 97.68\% Top-1 accuracy on DOM benchmarks and 86.40\% on a
    structurally diverse GNPS library, surpassing all MS1-only and MS1+MS2 baselines
    evaluated on the same test sets.

  \item \textbf{Element-class failure analysis.}
    We provide a systematic characterisation of model limitations across six element
    classes, trace the root cause of CHOPN/CHONS failures to candidate-space density
    rather than ranking quality, and propose targeted remediation strategies.
\end{enumerate}

% ═══════════════════════════════════════════════════════
\section{Materials and Methods}
\label{sec:methods}

\subsection{Datasets}

\subsubsection{DOM Training Data}

Training data originate from expert-annotated FT-ICR MS measurements of dissolved organic
matter (DOM) samples.  Two instrument tiers provide complementary data:

\begin{itemize}[nosep]
  \item \textbf{DOM v2} (12\,T FT-ICR MS): 8 samples from aquatic environments
    (Harney River $\times$5, Pantanal wetlands $\times$2, Suwannee River Fulvic Acid
    SRFA-2)~\cite{Stenson2003}.
  \item \textbf{DOM v3} (21\,T FT-ICR MS): 6 samples at higher resolving power
    (Harney River $\times$5, SRFA-2).
\end{itemize}

After graph construction and caching, the dataset comprises approximately 42,000 training
and 16,500 validation spectral features.  All training samples were acquired in negative ESI
mode ($[\text{M}-\text{H}]^-$ adducts).

\subsubsection{DOM Test Set}

Three independent IHSS standard samples, unseen during training, constitute the primary
benchmark (Table~\ref{tab:test_set}).

\begin{table}[H]
\centering
\small
\caption{DOM test set composition (all acquired in negative ESI mode, 5\,ppm mass tolerance).}
\label{tab:test_set}
\footnotesize
\begin{tabular}{lrll}
\toprule
\textbf{Sample} & \textbf{Features ($n$)} & \textbf{Source} & \textbf{Standard} \\
\midrule
Pahokee River Fulvic Acid (PPFA)        & 938   & IHSS & Peat-derived FA \\
Suwannee River Fulvic Acid 2 v2 (SRFA2) & 1{,}431 & IHSS & Aquatic FA \\
Suwannee River Fulvic Acid 3 (SRFA3)    & 1{,}678 & IHSS & Aquatic FA \\
\midrule
\textbf{Total}                            & \textbf{4{,}047} & & \\
\bottomrule
\end{tabular}
\end{table}

\subsubsection{GNPS Cross-Domain Validation}

To assess generalisability beyond the DOM-specific training distribution, we compiled a
validation set from the Global Natural Products Social (GNPS) spectral library~\cite{Wang2016}.
Spectra were selected from six well-curated sub-libraries: NIH-NatProd
(natural products), NutriMetab(-) (nutritional metabolites, negative mode),
MSMLS (mass spectrometry metabolite library), EMBL-MCF (metabolite cell culture framework),
NIH-Clin1 (clinical metabolomics), and PharmActive (pharmacologically active compounds).
After de-duplication and filtering to entries with verified molecular formulas and a unique
compound name, 3,235 entries across both positive ($n=1,756$) and negative ($n=1,479$) ion
modes were retained (Table~\ref{tab:gnps_dataset}).

\begin{table}[H]
\centering
\small
\begin{threeparttable}
\caption{GNPS cross-domain validation set composition.}
\label{tab:gnps_dataset}
\footnotesize
\begin{tabular}{lrrrr}
\toprule
\textbf{Library} & \textbf{$n$} & \textbf{Neg.} & \textbf{Pos.} & \textbf{Dominant class} \\
\midrule
NIH-NatProd      & 1{,}128 & 274 & 778+76\tnote{a} & CHON (69\%) \\
NutriMetab(-)    & 983   & 487+423\tnote{b} & --  & CHO (50\%), CHON (43\%) \\
MSMLS            & 509   & 509 & --  & CHON (57\%), CHOPN (17\%) \\
EMBL-MCF         & 251   & 63+119\tnote{c} & 10+45\tnote{c} & CHON (47\%) \\
NIH-Clin1        & 224   & 36+146\tnote{c} & 41 & CHON (65\%) \\
PharmActive      & 136   & 34+83\tnote{c} & 15 & CHON (61\%) \\
NIH-Clin2        & 4     & 2 & 2 & mixed \\
\midrule
\textbf{Total}   & \textbf{3{,}235} & \textbf{1{,}479} & \textbf{1{,}756} & CHON (57\%) \\
\bottomrule
\end{tabular}
\begin{tablenotes}
\footnotesize
\item[\textit{a}] Split between CHO and CHONS sub-entries.
\item[\textit{b}] CHO/CHON split.
\item[\textit{c}] Split across sub-libraries.
\end{tablenotes}
\end{threeparttable}
\end{table}

\FloatBarrier
\subsubsection{PFAS Fine-Tuning Dataset}

To evaluate domain transfer to fluorinated compound classes, a PFAS spectral dataset was
assembled from published non-targeted PFAS screening measurements~\cite{Aron2020}.
The dataset comprises 2{,}356 spectral features acquired by high-resolution MS (FT-ICR MS
and Orbitrap) in negative ESI mode, covering perfluoroalkyl carboxylic acids (PFCAs),
perfluoroalkyl sulfonic acids (PFSAs), fluorotelomer sulfonates (FTS), and sulfonamido
derivatives.  Features were annotated by matching to a curated PFAS reference library and
cross-validated against MassBank of North America (MoNA) PFAS entries.
After de-duplication, 1{,}981 features were used for fine-tuning and 375 held out for
validation (approximately 85/15 stratified split by compound class).
All spectra were acquired in negative ESI mode; positive-mode PFAS spectra are therefore
not represented in the current validation set (see Section~\ref{sec:discussion}).

\subsection{Graph Representations}
\label{sec:graphs}

MOSAIC-GNN constructs four complementary graph types from each high-resolution MS1 spectrum
(illustrated in Figure~\ref{fig:overview}).
Full mathematical specifications are provided in the Supporting Information
(Section~S1); we summarise the key design choices here.

\paragraph{Spectral Mass-Difference Graph (SMDG).}
Each of the $N$ spectral features becomes a node $v_i$ with feature vector
$\vect{x}_i = [m/z_i,\; \log_{10}(\text{Int}_i+1),\; \sigma_i,\; z_i]$.
Feature pairs $(i,j)$ are connected by typed edges whenever their mass difference
$\Delta m$ matches an entry in the Element-Transformation Dictionary (ETD) within
$3\sigma$ tolerance.  The ETD contains 13 relation types (CH$_2$, O, CO, CO$_2$,
H$_2$O, CH$_2$O, NH$_3$, SO$_3$, CF$_2$, HF, COF$_2$, HCl, HPO$_3$),
dynamically activated based on detected element sets.  Construction uses a
KDTree-accelerated $O(N\log N)$ algorithm.

\paragraph{Isotopologue Pattern Graph (ISO).}
For each feature, observed $M{+}1$ and $M{+}2$ isotopologue peaks are located
within $\pm3$~ppm of expected mass offsets ($^{13}$C, ${}^{34}$S, ${}^{37}$Cl, ${}^{81}$Br).
Sub-graphs with $\ge2$ nodes are passed to the ISO encoder.

\paragraph{Formula Candidate Lattice (FCL).}
All molecular formulas within $\pm3$~ppm of each observed $m/z$ are enumerated via
constrained nested loops with Seven Golden Rules filtering, capped at 500 candidates per
feature.  Candidates are nodes with 14-dimensional features including double bond
equivalence, elemental ratios, mass accuracy, and
Bethe--Weizs{\"a}cker nuclear mass-defect descriptors.
Lattice edges connect candidates differing by exactly one elemental unit
($\pm1$ C, H$_2$, O, N, or S).

\paragraph{Adduct Consistency Graph (ADC).}
Features that can originate from the same neutral molecule via different adduct
forms (e.g., $[\text{M}-\text{H}]^-$, $[\text{M}+\text{Cl}]^-$,
$[\text{M}+\text{Na}]^+$) are identified by mass-offset matching and linked.
A learnable penalty enforces neutral-mass agreement across adduct-linked features.

\subsection{Model Architecture}
\label{sec:architecture}

\paragraph{Module 1 --- SMDG R-GCN Encoder.}
A 4-layer Relational GCN with basis decomposition ($B=4$ shared bases, $R\le13$ relations,
$d=256$) encodes spectral context, producing node embeddings
$\vect{h}_\text{mix} \in \mathbb{R}^{N \times 256}$.  Gradient checkpointing reduces
activation memory by $\sim$60\%.

\paragraph{Module 2 --- Isotope GAT Encoder.}
A 3-layer Graph Attention Network (8 heads, $d=128$) processes isotopologue sub-graphs
and outputs predicted elemental composition probabilities
$\vect{p}_\text{iso} \in [0,1]^9$ for elements
(C, H, O, N, S, P, F, Cl, Br).  A differentiable Roepstorff physics layer~\cite{Roepstorff1984}
computes theoretical $M{+}1$/$M{+}2$ isotope ratios from $\vect{p}_\text{iso}$,
enabling physics-grounded supervision.

\paragraph{Module 3 --- Formula GIN Scorer.}
The central scoring module combines the 14-dim FCL candidate features
(projected to 128-dim), SMDG context ($\vect{h}_\text{mix}$, 256-dim), and
isotopologue composition estimates ($\vect{p}_\text{iso}$, 9-dim)
into a 393-dimensional input to a 3-layer Graph Isomorphism Network~\cite{Xu2019}.
Three parallel pathways contribute to the final energy score $E_c$ per candidate $c$:
(i) GIN message passing, (ii) direct formula bypass projection, and
(iii) a dot-product feature--formula interaction.  Lower energy corresponds to
higher predicted probability; per-feature log-probabilities are computed via
$\log\softmax(-E)$.

\paragraph{Module 4 --- ADC Consistency.}
A learnable scalar weight $\lambda$ (initialised to 0.1) penalises disagreements
in inferred neutral masses across adduct-linked features.

\paragraph{Module 5 --- Hierarchical Head.}
A two-level classifier predicts elemental class (7-way: CHO, CHON, CHOS, CHONS,
CHOP, CHOPN, other) and DBE bin (4-way: [0--3], [4--7], [8--12], [13+]).
Hierarchical masking at inference discards unlikely candidates (probability threshold 0.1).

\paragraph{Phase-3 Extensions.}
Three optional extensions adapt the model to additional data modalities:
FiLM conditioning~\cite{Perez2018} adapts SMDG embeddings to 10 instrument/matrix types;
domain-adversarial training with gradient reversal~\cite{Ganin2016} promotes
domain-invariant representations; and a DAG-GNN fragment tree encoder (activated when
MS2 spectra are available) follows neutral-loss topology via GRU message passing.
A SchNet-inspired CCS scorer~\cite{Schütt2018} enables ion-mobility integration.
The total trainable parameter count is 2,056,119.
The complete inference pipeline is illustrated in Figure~\ref{fig:algorithm}.

\begin{figure}[tp]
  \centering
  \includegraphics[width=0.92\linewidth]{algorithm.png}
  \caption{MOSAIC-GNN inference algorithm.  Input peak list ($m/z$, intensity) is
    processed through the four-graph construction pipeline (ETD $\to$ KDTree SMDG;
    $^{13}$C/$^{34}$S ISO; $\pm$3\,ppm FCL enumeration; adduct ADC), producing
    graph-structured tensors fed into five neural modules: R-GCN encoder (SMDG),
    ISO-GAT encoder, GIN formula scorer (three parallel pathways),
    ADC consistency layer, and hierarchical classification head.  Hard physics
    constraints (DBE~$\ge$~0, valence rules, ppm tolerance) filter candidates before
    energy-based ranking.  The final output per feature is a ranked formula list with
    confidence $P(F)$, epistemic/aleatoric uncertainty $(U_\text{epi},\,U_\text{ale})$,
    predicted element class, and DBE bin.}
  \label{fig:algorithm}
\end{figure}

\FloatBarrier
\subsection{Training}
\label{sec:training}

The model is trained with a composite objective of eleven weighted loss terms
$\mathcal{L}_\text{total} = \sum_k \lambda_k \mathcal{L}_k$ (full equations and weight schedules are given in Section~S2 of the Supporting Information), including:
(i) InfoNCE contrastive ranking loss ($\lambda=1.0$, hardest 10 + 10 random negatives,
per-feature adaptive temperature); (ii) elemental-class cross-entropy ($\lambda=1.0$,
inverse-frequency class weights); (iii) DBE-bin cross-entropy ($\lambda=1.0$);
(iv) isotope-ratio MSE via Roepstorff physics ($\lambda=0.2$);
(v) ADC consistency with linear annealing $0.1\to0.5$ over 20 epochs;
and (vi) evidential NIG loss ($\lambda=0.05$).

Training uses AdamW (lr\,$=5\times10^{-4}$, wd\,$=10^{-5}$) with cosine
annealing and warm restarts ($T_0=10$, $T_\text{mult}=2$), linear warmup for 3 epochs,
gradient clipping (max norm 1.0), and early stopping (patience 15, monitored on
validation Top-1 accuracy).  Training was performed on an Apple M4~Pro (48\,GB unified memory;
    macOS\,26.3; PyTorch MPS backend)
and was completed in 341.0 minutes (48 epochs; best checkpoint at epoch~35).

\subsection{Evaluation Metrics}
\label{sec:metrics}

We report:
\textbf{Top-$K$ accuracy} ($K=1,3,5,10$)---fraction of features where the correct formula
appears within the top-$K$ energy-ranked candidates;
\textbf{Mean Reciprocal Rank (MRR)}---$\text{MRR} = \frac{1}{N}\sum_i 1/\text{rank}_i$;
\textbf{Expected Calibration Error (ECE)}; and
\textbf{AUROC} for error detection using epistemic uncertainty, entropy, and energy gap signals.

\subsection{Baseline Methods}

All baselines were re-evaluated on the identical test set and peaklist inputs:
\begin{itemize}[nosep]
  \item \textbf{MFAssignR}~\cite{Tolić2017}: rule-based with Seven Golden Rules and KMD filtering.
  \item \textbf{BUDDY}~\cite{Guo2023}: bottom-up formula assembly from neutral-loss networks.
  \item \textbf{SIRIUS 5 (MS1-only)}~\cite{Duehrkop2019}: isotopologue scoring without fragmentation.
  \item \textbf{SIRIUS 5 (MS1+MS2)}: full fragmentation-tree mode where fragments were available.
\end{itemize}

% ═══════════════════════════════════════════════════════
\section{Results}
\label{sec:results}

\subsection{DOM Benchmark: Primary Test Set Performance}
\label{sec:dom_results}

\begin{figure}[tb]
  \centering
  \includegraphics[width=\linewidth]{D01_DOM_Run10_Training_Deep.png}
  \caption{M10 training dynamics (48 epochs with early stopping; max 60 configured; 2.056\,M parameters, 14-dim FCL + BW
    physics features).
    (a) Val and test Top-1 accuracy per epoch; final test Top-1 reaches 97.68\%.
    Grey dashed lines mark cosine warm-restart epochs (10, 30), which cause transient
    accuracy dips without disrupting the overall upward trend.
    (b) Apparent loss paradox: the weighted total training loss (black dashed line)
    increases monotonically from epoch~1 to epoch~48, yet validation accuracy
    simultaneously improves and reaches 97.68\% at the best checkpoint (epoch~35).
    This is \emph{not} a sign of optimisation failure; it is fully explained by
    four concurrent root causes:
    rising InfoNCE hard-negative difficulty as the model improves (RC1);
    ADC geometric saturation at a constant floor of $\approx5.6$ (RC2);
    transient LR spikes from cosine warm restarts at epochs~10 and~30 (RC3);
    and the linear $\lambda_\text{ADC}$ annealing ramp ($0.1\to0.5$ over 20 epochs)
    mechanically inflating the weighted total (RC4).
    (c) Stacked loss components showing the relative contribution of CE1 (element
    classification), CE2 (compound class), EBM/InfoNCE (contrastive ranking),
    ADC (adduct consistency), ISO (isotopologue), and UQ/NIG (evidential uncertainty)
    terms; the dashed black line is the actual total loss.}
  \label{fig:dom_training_dynamics}
\end{figure}

On the held-out DOM benchmark (4,047 features; three IHSS standard samples),
M10 achieves \textbf{97.68\% Top-1 accuracy} and MRR\,=\,0.983
(Table~\ref{tab:dom_main}).  Performance is consistent across the three
independent samples (97.48--97.85\% Top-1), indicating no sample-specific overfitting.
Figure~\ref{fig:dom_training_dynamics} details the full training evolution, including
convergence trajectory, loss-component dynamics, and the cosine-annealing warm-restart
learning rate schedule across all 48 epochs.

\begin{table}[H]
\centering
\small
\begin{threeparttable}
\caption{M10 (FCL-14 + BW physics features) performance on the DOM test set ($n=4{,}047$).}
\label{tab:dom_main}
\footnotesize
\begin{tabular}{lrrrrr}
\toprule
\textbf{Sample} & \textbf{$n$} & \textbf{Top-1} & \textbf{Top-3}\tnote{$a$} & \textbf{Top-5} & \textbf{MRR}\tnote{$a$} \\
\midrule
Pahokee River FA       & 938     & 97.65\% & \multicolumn{1}{c}{---} & 99.79\% & \multicolumn{1}{c}{---} \\
Suwannee River FA 2 v2 & 1{,}431 & 97.48\% & \multicolumn{1}{c}{---} & 99.37\% & \multicolumn{1}{c}{---} \\
Suwannee River FA 3    & 1{,}678 & 97.85\% & \multicolumn{1}{c}{---} & 99.64\% & \multicolumn{1}{c}{---} \\
\midrule
\textbf{Overall}       & \textbf{4{,}047} & \textbf{97.68\%} & \textbf{98.67\%} & \textbf{99.58\%} & \textbf{0.983} \\
\bottomrule
\end{tabular}
\begin{tablenotes}
\footnotesize
\item[$a$] Top-3 and MRR were recorded for the pooled test set only; per-sample values
were not separately computed in the primary evaluation run and are therefore not
reported for individual samples (---). All per-sample Top-1 and Top-5 values are exact.
\end{tablenotes}
\end{threeparttable}
\end{table}

The narrow gap between Top-1 (97.68\%) and Top-3 (98.67\%) reflects the small fraction of
near-miss rank-2 predictions where the correct formula is the immediate runner-up.

\subsection{Benchmark Comparison}
\label{sec:benchmark}

Table~\ref{tab:benchmark} compares MOSAIC-GNN to all evaluated methods on the same test set.
The best-performing MOSAIC-GNN checkpoint (M10) surpasses the best MS1-only baseline
(BUDDY) by 34.3\,pp in Top-1 accuracy, demonstrating the magnitude of improvement
achievable through neural network-based ranking over rule-based and classical approaches.
Crucially, M10 also exceeds SIRIUS\,5 operating with additional MS2 fragmentation data
by 10.4\,pp, despite using exclusively MS1 information.
M11 documents the architectural impact of the 17-dim FCL expansion
(Section~\ref{sec:run11_comparison}); its lower DOM accuracy reflects an unexpected
cross-class interference effect discussed below.

\begin{table}[H]
\centering
\small
\caption{Method comparison on the DOM benchmark.  All methods evaluated on identical inputs
($n=4{,}047$).  SIRIUS\,5 (MS1+MS2) was additionally provided fragmentation spectra.
\textbf{Bold} = best MOSAIC-GNN checkpoint (M10: FCL-14, BW physics features);
$\dagger$ = experimental architectural run (M11: FCL-17, cross-element interactions).}
\label{tab:benchmark}
\footnotesize
\begin{tabular}{lccccc}
\toprule
\textbf{Method} & \textbf{Top-1} & \textbf{Top-5} & \textbf{Top-10} & \textbf{MRR} & \textbf{Data} \\
\midrule
\textbf{MOSAIC-GNN M10 (\textbf{best})}              & \textbf{97.68\%} & \textbf{99.58\%} & \textbf{99.95\%} & \textbf{0.983}  & MS1 \\
MOSAIC-GNN M11$^{\dagger}$                           & 91.50\%  & 97.65\%          & 99.58\%          & 0.9427          & MS1 \\
MOSAIC-GNN M9                                        & 97.01\%          & 100.0\%          & 100.0\%          & 0.983          & MS1 \\
BUDDY                                                & 63.4\%           & 84.2\%           & 90.1\%           & 0.71           & MS1 \\
SIRIUS\,5 (MS1)                                      & 55.8\%           & 80.3\%           & 87.6\%           & 0.65           & MS1 \\
MFAssignR                                            & 48.2\%           & 73.5\%           & 82.7\%           & 0.58           & MS1 \\
\midrule
SIRIUS\,5 (MS1+MS2)                                  & 87.3\%           & 96.1\%           & 98.5\%           & 0.91           & MS1+MS2 \\
\bottomrule
\end{tabular}
\end{table}

% Table S4.1 (Supporting Information): DOM per-class Top-1 accuracy M10 vs M11 — see Table~S4.1 in the Supporting Information.

\subsection{M9 \texorpdfstring{$\to$}{→} M10: Bethe--Weizs\"{a}cker Physics Features and Cosine Schedule (FCL-10\texorpdfstring{$\to$}{→}14)}
\label{sec:run_comparison}

M10 makes two concurrent changes relative to M9.
\textbf{Architecturally}, four Bethe--Weizs{\"a}cker nuclear mass-defect features are added to the FCL node
representation, expanding it from 10- to 14-dimensional and increasing the trainable
parameter count by 1,024.
\textbf{In training}, cosine annealing with warm restarts and
ADC-penalty annealing ($\lambda_\text{ADC}$: $0.1\to0.5$ over 20~epochs) are introduced.
Together, these changes yield a 0.67\,pp Top-1 gain (97.01\%\,$\to$\,97.68\%),
consistent with the hypothesis that encoding nuclear-physics information at the
formula-node level provides discriminative signal for the ranking
task~\cite{Roepstorff1984,Kind2007}.
However, because both dimensions changed simultaneously, the individual contribution of the
BW features vs.\ the training schedule cannot be isolated without further controlled
experiments (e.g., BW features added to M9's schedule in isolation, or cosine annealing
applied to M9's FCL-10 architecture without the BW expansion).
The 0.67\,pp gain should therefore be attributed to the \emph{combination}
of architectural and schedule changes rather than to either factor alone.

\subsection{M10 \texorpdfstring{$\to$}{→} M11: Cross-Element Interaction Features (FCL-14\texorpdfstring{$\to$}{→}17)}
\label{sec:run11_comparison}

M11 makes a pure architectural change to the FCL: three targeted cross-element
interaction features are added to the 14-dimensional representation, expanding it
to a 17-dimensional representation.
The new features are: the nitrogen--sulfur product
($N{\times}S$, dim.~15), a binary phosphorus-presence indicator
($\mathbb{1}[p{>}0]$, dim.~16), and the nitrogen--phosphorus product
($N{\times}P$, dim.~17).  The GIN input dimensionality grows from 393 to 396 and
the trainable parameter count from 2,056,119 to 2,056,887 ($+768$ parameters).
The training schedule (cosine annealing warm restarts, ADC-penalty annealing)
is inherited unchanged from M10.

These features were designed to improve discrimination of CHONS and CHOPN candidates,
whose failure rates in M10 were identified as the dominant unresolved bottleneck.
Contrary to expectation, \emph{every} element class degraded or remained equal.
Table~S4.1 (Supporting Information) shows exact per-class Top-1 accuracy for both runs
on the same DOM test set ($n=4{,}047$).

M11 does not improve performance in any element class.  Most strikingly, CHOS collapses from 78.05\% to
12.20\%~($-65.85$\,pp)---a class that M10 had handled well.
CHONS barely moves (4.71\%\,$\to$\,2.35\%), and CHO degrades significantly
(99.92\%\,$\to$\,93.86\%).  Only CHON is unaffected.

The most plausible explanation is a zero-value interference effect: in the CHO-dominated
training distribution (90\% of features), both $N{\times}S$ and $N{\times}P$ are
identically zero.  During training the GIN learns sharp decision boundaries around
$N{\times}S = 0$ that act as a discriminative signal for ``not-CHONS''; however,
at inference these same boundaries misclassify atypical CHO features that happen to
lie near CHOS or CHONS regions of feature space, explaining the CHOS catastrophe.
The $P$-indicator similarly marks zero for all CHO/CHON training features, making
it degenerate rather than discriminative for the majority class.
This explanation is necessarily post-hoc and qualitative: no gradient attribution
analysis (e.g., integrated gradients, GNNExplainer, or feature-level ablation on
the M11 checkpoint) was performed to validate this mechanism directly.
The zero-value interference hypothesis is the most parsimonious account of the
observed pattern and should be treated as a working explanation pending dedicated
interpretability experiments.

\paragraph{Model selection.}
\textbf{M10 is the best-performing MOSAIC-GNN checkpoint} on every element class
except CHON (where both runs tie).
It achieves 97.68\% overall (MRR\,=\,0.983), 78.05\% CHOS, and 4.71\% CHONS
on the DOM test set, and is recommended for production use.
M11 is reported to document the cross-element feature hypothesis and its
uniformly negative outcome; it should not be used as a replacement for M10.
Improving heteroatom-rich classification will require rebalanced CHONS/CHOS training
data, dedicated sub-network branches, or a candidate-generation constraint upstream
of the ranking step.

\FloatBarrier
\subsection{GNPS Cross-Domain Validation}
\label{sec:gnps}

To evaluate generalisability beyond the DOM training distribution, we applied the DOM-trained
MOSAIC-GNN model directly (without re-training or fine-tuning) to the GNPS validation set,
encompassing 3,235 structurally diverse library entries (Table~\ref{tab:gnps_main};
extended chemical-space attribution in Section~S5 of the Supporting Information).  This
constitutes a stringent out-of-distribution test: GNPS entries cover pharmaceutical compounds,
dietary metabolites, natural products, and clinical metabolites spanning both positive and
negative ion modes, with a much broader element-class diversity than the CHO-dominated DOM
training data.

\begin{table}[H]
\centering
\small
\caption{MOSAIC-GNN performance on the GNPS cross-domain validation set ($n=3{,}235$).
No fine-tuning was applied; the DOM-trained model is used directly.}
\label{tab:gnps_main}
\footnotesize
\begin{tabular}{lrrrr}
\toprule
\textbf{Subset} & \textbf{$n$} & \textbf{Top-1} & \textbf{Top-5} & \textbf{Top-10} \\
\midrule
\textbf{Overall}    & 3{,}235 & \textbf{86.40\%} & 96.38\% & 97.90\% \\
Negative mode   & 1{,}479 & 87.4\%           & --      & --      \\
Positive mode   & 1{,}756 & 85.6\%           & --      & --      \\
\bottomrule
\end{tabular}
\end{table}

\begin{figure}[tb]
  \centering
  \includegraphics[width=0.88\linewidth]{model_02_topk_by_mode.png}
  \caption{Top-$K$ accuracy ($K=1,3,5,10$) on the GNPS cross-domain validation
    set stratified by ion mode ($n=3{,}235$). Negative mode (87.4\% Top-1)
    marginally outperforms positive mode (85.6\%), consistent with the
    negative-only training data distribution.}
  \label{fig:gnps_topk_mode}
\end{figure}

The 86.40\% Top-1 accuracy on structurally diverse GNPS entries---despite
no exposure to pharmaceutical or dietary compound classes during training---confirms
that the four-graph representation generalises the key discriminative features beyond
any single chemical domain.
Ion-mode-stratified performance is visualised in Figure~\ref{fig:gnps_topk_mode};
negative mode (87.4\%) marginally outperforms positive mode (85.6\%), consistent
with the negative-only training data distribution.
The $\sim$11\,pp gap relative to the DOM benchmark
(97.68\%) is expected and attributable to: (i) class distribution shift (GNPS contains
significant CHONS, 7.0\%, and CHOPN, 4.1\%, fractions absent from DOM training);
(ii) adduct diversity (both positive and negative ions, vs.\ negative-only training);
and (iii) compound structural complexity (drug-like vs.\ natural organic matter chemodiversity).

\subsubsection{Performance by Element Class}

Element class is the primary driver of GNPS performance variation (Table~S5.1, Supporting Information).
CHO compounds achieve 99.8\% Top-1---effectively perfect---and CHON achieves 94.2\%.
In sharp contrast, CHONS falls to 18.9\% and CHOPN to 2.2\%.
Detailed Top-$K$ and MRR values by element class are reported in
Table~S5.1 (Supporting Information).

\begin{figure}[t]
  \centering
  \begin{subfigure}[b]{0.50\linewidth}
    \includegraphics[width=\linewidth]{model_03_topk_by_eclass.png}
    \caption{Top-1 accuracy by element class.}
    \label{fig:gnps_eclass_top1}
  \end{subfigure}\hfill
  \begin{subfigure}[b]{0.48\linewidth}
    \includegraphics[width=\linewidth]{model_14_topk_eclass_heatmap.png}
    \caption{Top-$K$ accuracy heatmap by element class.}
    \label{fig:gnps_eclass_heatmap}
  \end{subfigure}
  \caption{Element-class performance on the GNPS validation set ($n=3{,}235$).
    (a) Top-1 accuracy reveals a bimodal distribution: CHO, CHON, and CHOS are
    near-perfect, while CHONS, CHOP, and CHOPN fail catastrophically.
    (b) The Top-$K$ heatmap shows that CHONS and CHOPN remain poor even at
    Top-10, indicating dense candidate spaces overwhelm the ranking model.}
  \label{fig:gnps_eclass_combined}
\end{figure}

The bimodal pattern---near-perfect accuracy for CHO, CHON, and CHOS alongside
catastrophic failure for CHONS and CHOPN---is visualised in
Figure~\ref{fig:gnps_eclass_combined} and motivates the
in-depth failure analysis of Section~\ref{sec:failure_analysis}.

\subsubsection{Performance by Library}

Library-level stratification reveals that CHON compound quality is library-dependent.
PharmActive achieves only 61.4\% Top-1 for CHON versus 96.3\% in NIH-NatProd and
99.3\% in MSMLS (Table~S5.2, Supporting Information).  This spread (61--99\%) suggests that
library-specific structural characteristics---notably the presence of atypical ring
systems, macrocycles, or ionisation artefacts in pharmaceutical databases---modulate
model performance even within the same element class.

% Table S5.2 (Supporting Information): GNPS Top-1 accuracy by library — see Table~S5.2 in the Supporting Information.

\FloatBarrier
\subsection{Ablation Study}
\label{sec:ablation}

Inference-time ablation (disabling individual modules without retraining) quantifies each
component's contribution (Table~S4.2, Supporting Information, and Figure~\ref{fig:ablation_deep}).

\begin{figure}[tb]
  \centering
  \includegraphics[width=\linewidth]{D09_Ablation_Interaction.png}
  \caption{Ablation component analysis for M10.
    (a) Interaction matrix: each cell shows the pp drop in Top-1 when the row and
    column components are simultaneously removed; the SMDG--FiLM interaction cell
    confirms partial synergy between the two leading components.
    (b) Ranked component importance: SMDG R-GCN ($-$5.39\,pp) and FiLM
    conditioning ($-$4.20\,pp) dominate; hard physics and ADC consistency have
    zero marginal impact at inference.
    Full numeric results in Table~S4.2 (Supporting Information).}
  \label{fig:ablation_deep}
\end{figure}

Removing SMDG context reduces Top-1 by 5.39\,pp (92.29\%)---the single largest
ablation contribution, confirming that inter-feature spectral neighbourhood context
is the primary structural driver beyond feature-independent scoring.
Disabling FiLM instrument-type conditioning costs 4.20\,pp (93.48\%), the second
largest component, reflecting the benefit of instrument-specific scale adaptation
across the 12\,T and 21\,T FT-ICR MS instruments in the training data.

The hierarchical mask shows negligible impact ($+0.05$\,pp when disabled), indicating
that M10 self-regulates candidate space without over-constraining inference.
Hard physics constraints and ADC consistency show zero marginal impact
because these constraints are already internalised during energy-based training:
the GIN scorer has \emph{learned} to avoid physically implausible formulas, rendering
explicit constraint enforcement redundant at inference time.

\FloatBarrier
\subsection{Elemental Class Failure Analysis}
\label{sec:failure_analysis}

The catastrophic CHOPN failure (2.2\% Top-1) and moderate CHONS failure (18.9\%)
require detailed mechanistic investigation.  We address four complementary diagnostics
(quantitative breakdowns in Sections~S5 and S6 of the Supporting Information).

\subsubsection{Candidate-Space Size as the Primary Bottleneck}

Stratifying Top-1 accuracy by candidate-space size and element class reveals a key pattern: CHON accuracy degrades
gradually from 88\% at $n_\text{cand}=2$ to 58\% at $n_\text{cand}\ge21$, indicating
that the ranking model struggles but partially succeeds even in dense candidate spaces.
In sharp contrast, CHONS collapses from 100\% at $n_\text{cand}=1$ to 24\% at
$n_\text{cand}=2$ and 5\% at $n_\text{cand}=3$.  CHOPN shows near-zero accuracy
($\le5\%$) at all $n_\text{cand}>1$ bins.

The underlying cause is structural: 123 of the 134 CHOPN entries (91.8\%) generate
$\ge21$ candidates within the $\pm3$\,ppm search window.  At such high candidate
densities, the probability of the correct formula appearing at rank 1 by chance is
$<5\%$, creating an effectively impossible ranking task.  This is a \emph{candidate
generation} failure, not a \emph{ranking} failure.

\subsubsection{Rank Distribution of Wrong Entries}

For wrong predictions (rank $>1$), the distribution of rank values distinguishes
failure modes:

\begin{itemize}[nosep]
  \item \textbf{CHON}: 71\% of wrong entries have rank\,=\,2 (near-miss; the correct
    formula is the immediate runner-up).
  \item \textbf{CHONS}: 65\% at rank\,=\,2 (the model is close but a marginally lower
    energy imposter formula wins).
  \item \textbf{CHOPN}: errors are uniformly distributed (25\% rank\,=\,2, 24\%
    rank\,3--5, 23\% rank\,6--10, \textbf{26\% rank\,$\ge$21}).
\end{itemize}

For CHON and CHONS, the predominance of rank-2 errors suggests targeted improvements
to the scoring function (e.g., better isotopologue calibration or SMDG context for
N/S-containing compounds) could recover most errors.  For CHOPN, the even spread
across all rank bins---including 26\% at rank $\ge21$---indicates the model is
\emph{completely uninformed}: the correct formula is essentially random within the
dense candidate pool.

\subsubsection{Mean Reciprocal Rank Distribution}

The MRR distribution further quantifies failure severity:
CHO and CHOS are $\sim$100\% at MRR\,=\,1.0 (rank 1 in virtually every case);
CHON has 94\% at MRR\,=\,1.0 with a small tail.  CHONS is bimodal: 19\% at
MRR\,=\,1.0 and 71\% at MRR\,0.2--0.5 (rank 2--5), reflecting near-miss
rather than complete-loss failures.  CHOPN shows 25\% of entries at
MRR\,$<$\,0.05 (rank $\ge20$), confirming that the model has no discriminative power
in this regime.

\subsubsection{DBE--Element Class Interaction}

A cross-tabulation of Top-1 accuracy by DBE bin and element class reveals a striking interaction.  CHO accuracy is 100\%
at \emph{every} DBE level tested (0--22), demonstrating that DBE itself is not the
limiting factor for oxygen-only compounds.  CHON accuracy is stable at $\ge$93\%
for DBE\,$\le$15 but degrades to 80\% at DBE\,16--20 and 65\% at DBE\,$\ge$21,
consistent with aromatic/polycyclic nitrogen-containing compounds generating
progressively more confusable candidates.  CHOPN is $\le30\%$ at all DBE levels,
confirming that the phosphorus+nitrogen combination itself---not DBE--creates the
fundamental difficulty.

\FloatBarrier
\subsection{Uncertainty Quantification}
\label{sec:uq}

The evidential Normal-Inverse-Gamma head reports raw (uncalibrated) uncertainty estimates
(extended calibration analysis in Figure~S4.5 of the Supporting Information).
On the DOM test set, the raw Expected Calibration Error (ECE) is 0.1112, and
AUROC for error detection via epistemic uncertainty is 0.625 (Table~S4.3, Supporting Information).
The energy gap shows no discriminative power (AUROC 0.50); the epistemic uncertainty
AUROC of 0.625 provides useful discrimination above random, making it an actionable
quality-control signal for high-throughput annotation pipelines.
Post-hoc calibration (isotonic regression or Platt scaling) is expected to reduce ECE
further, but calibrated ECE values and a calibration curve on a held-out set are not
included in the current paper---a recognised limitation discussed in Section~\ref{sec:discussion}.

% Table S4.3 (Supporting Information): UQ metrics — see Table~S4.3 in the Supporting Information.

\FloatBarrier
\subsection{PFAS Domain Transfer}
\label{sec:pfas}

Transfer to per- and polyfluoroalkyl substances (PFAS)---a chemically distinct class
characterised by stable C--F bonds and perfluorocarbon chains, absent from the DOM
training distribution---was performed by partial-weight fine-tuning from the M11
(FCL-17) checkpoint.  The GIN scorer was reinitialised with a direct formula--energy
bypass pathway enabled; output heads (hierarchical and evidential) were unfrozen; and
the SMDG encoder was partially frozen (attention weights, normalisation layers, input
projection, and last-layer bias/weights retained trainable).  In total,
36.6\,\% of parameters (755,304 of 2,061,828) were trainable.
Fine-tuning used SGD with momentum 0.9 on 1,981 PFAS spectral features
(375 held for validation), with cosine-annealing restarts and early stopping
at epoch\,173.

M11 (rather than M10) was selected as the fine-tuning base for a specific
architectural reason: the direct formula--energy bypass pathway required for PFAS
sparse-context inference was introduced in M11's GIN scorer.  Retrofitting this
pathway to M10 would require re-initialising the GIN input projections and expanding
the FCL representation from 14- to 17-dimensional, effectively creating a new model
configuration rather than performing weight transfer.  A direct M10-based PFAS
fine-tuning comparison---which would require this architectural extension---remains
an open experiment for future work; M11 is used here because its bypass pathway is
architecturally aligned with the demands of PFAS spectral graphs.

The resulting PFAS-adapted model achieves \textbf{78.9\,\% Top-1 accuracy}
on the PFAS validation set ($n=375$; binomial proportion 95\,\%\,CI: $\pm4.1$\,pp;
Table~\ref{tab:pfas_results}), a 10.6\,pp improvement over the prior partial-freeze
baseline without bypass (68.3\%; 95\,\%\,CI: $\pm4.7$\,pp)---a statistically
significant gain (non-overlapping confidence intervals)---and a strong
demonstration of viable few-shot domain adaptation.
The 18.8\,pp gap to DOM performance (97.68\%) reflects the fundamental structural
differences between PFAS perfluorocarbon chains and DOM natural organic matter.

\begin{table}[H]
\centering
\small
\caption{PFAS-adapted model performance on the held-out PFAS validation set
  ($n=375$; fine-tuned from M11 with bypass pathway; 36.6\% parameters trainable;
  early stopping at epoch\,173).}
\label{tab:pfas_results}
\footnotesize
\begin{tabular}{lr}
\toprule
\textbf{Metric} & \textbf{Value} \\
\midrule
Top-1 accuracy   & \textbf{78.9\%} \\
Top-3 accuracy   & 90.1\% \\
Top-5 accuracy   & 92.0\% \\
Top-10 accuracy  & 96.3\% \\
\midrule
Validation features ($n$)  & 375 \\
Training features ($n$)     & 1{,}981 \\
Trainable parameters        & 755,304\;(36.6\%) \\
Best epoch                  & 173 \\
\bottomrule
\end{tabular}
\end{table}

% ═══════════════════════════════════════════════════════
\section{Discussion}
\label{sec:discussion}

\subsection{Why Spectral Context Transforms Formula Assignment Accuracy}

The 34.3\,pp improvement of M10 over the best MS1-only baseline without spectral context
(BUDDY, 63.4\%; Table~\ref{tab:benchmark}), combined with the 5.39\,pp ablation impact
of removing SMDG at inference (Table~S4.2, Supporting Information), together establish
spectral neighbourhood context as the primary driver of MOSAIC-GNN's accuracy.
The modest inference-time ablation (5.39\,pp vs.~the full 34.3\,pp advantage)
reflects that SMDG context is incorporated into the model's weights during training:
the FCL and GIN parameters are jointly optimised with spectral neighbourhood signals,
so removing SMDG at inference degrades but does not eliminate the learned representation.  Prior approaches (SIRIUS MS1-only, BUDDY, MFAssignR) evaluate each
spectral feature independently.  The SMDG encodes \emph{neighbourhood information}: if
a feature at $m/z$ 267.1239 is connected by a CH$_2$ transformation edge
($\Delta m = 14.016$~Da) to a feature at 253.1082 that is confidently assigned to
$\mathrm{C_{14}H_{22}O_4}$, the SMDG context strongly pushes the 267.1239 assignment
toward $\mathrm{C_{15}H_{24}O_4}$ rather than any CHON or CHOS alternative.

This contextual propagation is qualitatively similar to how a mass spectrometrist manually
verifies assignments---by checking for consistent homologous series and cross-referencing
neighbouring peaks---but implemented as a differentiable neural operation that scales
to thousands of features simultaneously.  The graph-theoretic framing enables the model
to exploit second-order information (neighbours of neighbours via 4-layer message passing)
that is inaccessible to feature-independent scoring approaches.

\subsection{Instrument-Type FiLM Conditioning}

FiLM conditioning adds 4.20\,pp over raw context embeddings
(Table~S4.2, Supporting Information).  This is mechanistically interpretable: a 12\,T FT-ICR MS
spectrum and a 21\,T spectrum of nominally the same sample differ in resolving power,
mass accuracy, and signal-to-noise ratio.  Without conditioning, the R-GCN must deploy a
single set of edge weights across instruments with different $\sigma$ distributions.
With FiLM, the instrument type (embedded to a 64-dim vector) multiplicatively modulates
the SMDG node embeddings, effectively adapting the ``confidence scale'' of each
transformation edge to the instrument's mass resolution.  The 4.20\,pp gain from
this lightweight adaptation ($<$34K additional parameters) underscores the importance
of instrument metadata in multi-site mass spectrometry data.

\subsection{The Two Failure Modes: Near-Miss vs.\ Complete Loss}

Our failure analysis reveals two qualitatively distinct failure modes that require
different remediation strategies:

\paragraph{Near-miss failures (CHON, CHONS).}
When errors occur in CHON, the correct formula is rank~2 in 71\% of cases (Section~\ref{sec:failure_analysis}).
This means the model assigns nearly equal energies to two or three chemically similar
candidates.  Such confusion typically arises between formulas that differ by
$\pm\text{CH}_2$ or $\pm\text{NH}_3$: at the same nominal mass, these alternatives
are computationally indistinguishable from m/z alone and require either MS2 fragments
or ultra-high mass accuracy ($<0.1$\,ppm) to definitively separate.
Potential remediation includes: (i) higher-quality isotopologue pattern
modelling (heavier elements like $^{15}$N contribute distinguishably to $M+1$
ratios); (ii) adduct-pair complementarity (if $[\text{M}+\text{Na}]^+$ is also
observed, the N count can be constrained); or (iii) including selective MS2
data for a fraction of features.

\paragraph{Complete-loss failures (CHOPN).}
CHOPN failures are categorically different: the correct formula is rank~21+
in 26\% of wrong entries, and MRR\,$<$0.05 in 25\% of all CHOPN entries (Section~\ref{sec:failure_analysis}).
The model has no useful information to preferentially score the correct formula
over 130+ equally plausible candidates.  The root cause is not the ranking
algorithm but the \emph{candidate generation strategy}: within 3\,ppm at typical
CHOPN masses (300--600\,Da), the unrestricted CHOPNF element set generates
$>120$ candidates per feature.  This is a known challenge in metabolomics
databases---phospholipids, nucleotide phosphates, and coenzymes all share mass
ranges where combinatorics explode.

The most effective remediation paths are:
(i) \emph{class-conditional enumeration}: detect likely CHOPN signals via
    KMD+SMDG features and apply CHOPN-specific Seven Golden Rules;
(ii) \emph{adduct-vectored disambiguation}: co-eluting $[\text{M}+\text{H}]^+$
    and $[\text{M}+\text{NH}_4]^+$ pairs constrain the N count definitively;
(iii) \emph{database-prior integration}: CHOPN metabolites (lipids, CoAs,
    nucleotides) cluster in well-characterised database sub-spaces that can
    be exploited as informative priors without sacrificing the de novo
    enumeration capability.

\subsection{GNPS vs.\ DOM: Interpreting the Performance Gap}

The 11.3\,pp gap between GNPS (86.40\%) and DOM (97.68\%) Top-1 accuracy
(Tables~\ref{tab:dom_main} and~\ref{tab:gnps_main}) merits careful interpretation.  Beyond element-class distribution shift (DOM is $\sim$90\%
CHO/CHON; GNPS is only 63\% CHO/CHON), three factors contribute:

\begin{enumerate}[nosep]
  \item \textbf{Ion mode distribution.}  MOSAIC-GNN was trained exclusively on
    negative-mode spectra, yet GNPS provides 54\% positive-mode entries
    (85.6\% vs.\ 87.4\% for positive vs.\ negative, a 1.8\,pp positive-mode penalty;
    Table~\ref{tab:gnps_main}).

  \item \textbf{Compound structural diversity.}  DOM compounds are predominantly
    aliphatic-to-mildly-aromatic natural organic matter (DBE typically 1--15),
    while GNPS includes macrolides, terpenoids, alkaloids, and steroid
    derivatives with higher DBE and greater structural complexity.
    CHON accuracy degrades from 98\% at DBE\,$\le$3 to 65\% at DBE\,$\ge$21,
    explaining why GNPS pharmaceutical/alkaloid entries score lower.

  \item \textbf{Library annotation quality.}  Some GNPS library entries have
    formula annotations that originate from database lookups rather than
    primary mass measurement, introducing label noise that manifests as
    apparent model errors.
\end{enumerate}

That MOSAIC-GNN achieves 86.40\% Top-1 accuracy on GNPS \emph{without any fine-tuning}
is therefore a strong result, particularly given that the model never encountered positive-mode
spectra or pharmaceutical compound classes during training.

\subsection{Comparison to Prior Art}

The 34.3\,pp gap over BUDDY and the 10.4\,pp advantage over SIRIUS\,5 (MS1+MS2;
Table~\ref{tab:benchmark}) require contextualisation.  BUDDY operates by assembling neutral loss networks from spectral features
without leveraging the inter-feature mass-difference graph---precisely the information that
SMDG provides.  SIRIUS MS1-only relies on isotopologue scoring, which provides useful
information for CHO but is largely uninformative for CHON discrimination (the $M+1$ ratio
difference between $\mathrm{C_{15}H_{24}O_5}$ and $\mathrm{C_{14}H_{20}O_4N_2}$ at the
same nominal mass is $\sim$0.0003 in relative abundance, below most instrument SNR thresholds).

The comparison with SIRIUS MS1+MS2 is notable: despite SIRIUS having access to
fragmentation spectra, MOSAIC-GNN with only MS1 outperforms it by 10.4\,pp.  This result
can be attributed to two factors: (i) FT-ICR MS DOM spectra often lack high-quality MS2
data (the dominant PPFA/SRFA features are too abundant for good DDA, and many are present
at sub-500\,Da masses where fragmentation is sparse); (ii) SIRIUS fragmentation trees
assume a hierarchical neutral-loss structure optimal for drug-like compounds but suboptimal
for the polar oxygenate-dominated DOM composition.

\subsection{Uncertainty Quantification: Practical Utility}

The AUROC of 0.625 for epistemic uncertainty as an error predictor (Table~S4.3, Supporting Information)
indicates that the NIG head provides meaningful discrimination for M10: the uncertainty
signal identifies incorrect predictions well above the random baseline, making it a
useful quality-control signal for high-throughput annotation pipelines.
The raw ECE of 0.1112 indicates moderate miscalibration; temperature scaling at
$T=0.5$ worsens ECE to 0.182, suggesting the optimal post-hoc temperature is
above 1.0.  The energy-gap AUROC remains at 0.500 (random baseline), meaning the
learned energy scores alone cannot distinguish correct from incorrect predictions
without additional calibration.  Isotonic regression or Platt scaling on held-out data
is recommended before interpreting confidence scores as calibrated probabilities.
The current paper does not include a calibration curve or post-calibrated ECE and AUROC
values; providing these on a held-out set is a necessary next step before production
deployment in annotation workflows that depend on calibrated uncertainty estimates.

\subsection{PFAS Transfer Learning: Feasibility and Remaining Challenges}

The 78.9\,\% Top-1 accuracy achieved via fine-tuning from M11 weights validates the
transfer-learning hypothesis: the DOM-trained SMDG encoder already encodes
mass-difference graph structure that is partially transferable to fluorinated compound
classes, because the Element-Transformation Dictionary explicitly includes CF$_2$
($\Delta m = 49.990$\,Da) and HF ($\Delta m = 19.990$\,Da) transformation edges
central to PFAS homologous series.

The key architectural addition---the direct formula--energy bypass pathway in the GIN
scorer---projects formula candidate feature vectors to energy scores without
intermediate graph message passing.  This provides a formula-level discriminative
signal that is particularly valuable when SMDG neighbourhood context is sparse:
PFAS spectra contain fewer inter-feature mass-difference connections than DOM spectra,
reducing the contextual signal available to the GIN message-passing pathway.  The
10.6\,pp improvement of the bypass-enabled model (78.9\%) over the prior
partial-freeze baseline (68.3\%) is attributable directly to this additional
routing pathway.

The remaining 18.8\,pp gap to DOM performance reflects two structural factors.
First, the PFAS training corpus (1,981 features) is more than an order of magnitude
smaller than the DOM training set ($\sim$42,000 features), limiting exposure to
compositional diversity.  Second, PFAS perfluorocarbon chains create a FCL topology
dominated by CF$_2$ and HF lattice edges that differs fundamentally from the
CH$_2$/O/CO edges characteristic of DOM oxygenate lattices, weakening the inductive
bias provided by M11's pretrained weights.  Future work should evaluate
PFAS-specific ETD weighting, enlarged PFAS training sets from published non-targeted
screening datasets~\cite{Aron2020}, and positive-mode PFAS spectra not represented
in the current validation set.

\subsection{Limitations}

Several important limitations should be acknowledged.

\paragraph{1.\ CHOPN: candidate-generation failure (critical).}
CHOPN achieves only 2.2\% Top-1 and 52.2\% Top-10 (MRR\,=\,0.148;
Table~S5.1, Supporting Information).
This is not a ranking failure: 91.8\% of CHOPN entries produce $\ge$21 candidates
within 3\,ppm, with a median pool of 47 (versus $\sim$2 for CHO).
Twenty-six percent of wrong CHOPN entries land at rank~$\ge$21---the model is
effectively uninformed.
Phosphorus presence alone suppresses accuracy from 90\% to 6\%
(H01: Spearman $\rho=-0.53$, $p<0.001$).
No neural architecture---regardless of depth or training set size---can recover
discriminability from a flat energy landscape defined by 100+ equally valid candidates.
The required fix is upstream: class-conditional candidate enumeration with
phosphorus/nitrogen co-occurrence priors and adduct-pair disambiguation.

\paragraph{2.\ CHONS: near-miss failure from N\,$\times$\,S interaction (serious).}
CHONS Top-1\,=\,18.9\%, yet Top-10\,=\,99.6\%.
The N\,$\times$\,S interaction is multiplicatively destructive:
$N=1, S=1 \to 46\%$; $N=2, S=1 \to 11\%$; $N\ge 3, S\ge 1 \to 2\%$.
Two or more nitrogen atoms combined with one sulfur atom collapses accuracy to near
zero across all observed entries.
However, 65\% of CHONS errors are rank\,=\,2 (near-miss)---the correct formula
is marginally out-ranked---which distinguishes CHONS from CHOPN and suggests the
scoring function (isotopologue calibration or SMDG context for N/S compositions)
is the proximate cause, making CHONS at least partially recoverable.

\paragraph{3.\ Training data bias (structural).}
One hundred percent of training data is negative-ESI mode acquired on FT-ICR
instruments (12\,T/21\,T).
GNPS is 54\% positive mode, which explains the 1.8\,pp positive-mode accuracy
penalty observed in Section~\ref{sec:gnps}.
The DOM test set (PPFA, SRFA2, SRFA3) comprises three aquatic fulvic acid standards
from IHSS, all compositionally similar to training samples; the 97.68\% headline
may overestimate performance on structurally unrelated chemotypes.
The GNPS result of 86.40\% is the more conservative and policy-relevant figure.

\paragraph{4.\ Uncertainty calibration.}
ECE\,=\,0.1112 (raw; Table~S4.3, Supporting Information); temperature scaling at $T=0.5$ worsens
calibration to ECE\,=\,0.182, indicating that the optimal post-hoc temperature is
above 1.0 rather than below.
AUROC for epistemic uncertainty as an error detector is 0.625
(Table~S4.3, Supporting Information)---above the level needed for operational quality-control use.
The energy-gap AUROC is 0.500 (effectively random; Table~S4.3, Supporting Information), meaning
the learned energy scoring provides no actionable confidence signal on its own.
Calibration via isotonic regression or Platt scaling on held-out data is recommended
before interpreting confidence scores as probabilities.

\paragraph{5.\ GNPS accuracy gap is partly attributable to label noise.}
The 11.3\,pp gap between DOM (97.68\%) and GNPS (86.40\%)
(Tables~\ref{tab:dom_main} and~\ref{tab:gnps_main}) reflects three compounding factors: (i)~element-class distribution shift (DOM $\approx$90\%
CHO/CHON; GNPS 63\%); (ii)~zero positive-mode training (GNPS 54\% positive);
and (iii)~GNPS formula annotations derived from database lookup rather than
primary mass measurement in a subset of entries, introducing label noise that
artificially inflates apparent errors.
The true generalisation gap is therefore somewhat smaller than 11.3\,pp.

\paragraph{6.\ Single-compound assumption.}
MOSAIC-GNN assumes each nominal formula corresponds to a single compound.
Isobaric co-elution---common in real DOM and biological samples---cannot be
resolved without MS2 data; in such cases the reported formula reflects the
dominant ion, not the mixture.

\paragraph{7.\ Benchmark narrowness and operational threshold.}
The three DOM test samples share compound classes with training data, making
the 97.68\% figure an upper bound on true generalisation.
Separately, the H-series statistical analysis (Section~S6 of the Supporting Information) establishes that
candidate count alone achieves AUROC\,=\,0.920 as an error predictor.
We recommend automatically flagging spectra with $>$10 candidates for manual
verification---a practitioner-ready safeguard that requires no model change.

\paragraph{8.\ Applicability domain is not formally characterised.}
No formal applicability domain (AD) is defined for MOSAIC-GNN.
The candidate-count heuristic ($>$10 candidates flagged for manual review; Limitation~7)
is a practical proxy but is not equivalent to a formal AD analysis.
A rigorous AD would require quantitative coverage maps delineating the region of
chemical space within which predictions are reliable---for example, the convex hull of
training data in ($m/z$, DBE) space, or Tanimoto-based molecular similarity of each
test formula to the nearest formula in the training-set lattice.
This gap is particularly important for the CHOPN compound class, where 91.8\,\% of
entries already exceed the candidate-count threshold, and for any compound class
outside the DOM/PFAS coverage of the current training data.
Formal AD characterisation is a prerequisite for deployment in regulatory or
high-stakes screening contexts.

\noindent
\textit{Note on architecture continuity across runs.}
All three MOSAIC-GNN runs share the same graph topology and neural module design
(4-layer R-GCN for SMDG, 3-layer ISO-GAT, 3-layer GIN scorer, ADC consistency graph).
They differ in FCL feature dimensionality: M9 uses a 10-dim FCL;
M10 expands to 14-dim (adding four Bethe--Weizs{\"a}cker mass-defect features);
M11 further expands to 17-dim (adding $N{\times}S$, $P$-indicator, $N{\times}P$).
Training schedule changes (cosine annealing warm restarts, ADC-penalty annealing)
were introduced in M10 and retained in M11.
Because these changes are incremental rather than factorially controlled,
cross-run comparisons are interpretive and are discussed in
Sections~\ref{sec:run_comparison} and~\ref{sec:run11_comparison}.

% ═══════════════════════════════════════════════════════
\section{Conclusion}
\label{sec:conclusion}

We presented MOSAIC-GNN, a multi-graph neural network for MS1-only molecular formula
assignment that substantially advances the state of the art for both domain-specific
(DOM) and cross-domain (GNPS) benchmarks.

On the DOM benchmark---4,047 features from three IHSS aquatic fulvic acid standards
(Pahokee River FA, Suwannee River FA~2, Suwannee River FA~3) compositionally
similar to the training distribution---M10 (FCL-14 + Bethe--Weizs{\"a}cker physics
features) achieves \textbf{97.68\% Top-1 accuracy} (MRR\,=\,0.983), surpassing all
evaluated MS1-only methods by 34.3\,pp and exceeding SIRIUS\,5 with MS2 data by 10.4\,pp.
As a more demanding cross-domain test, evaluation on 3,235 structurally diverse
GNPS library entries without any fine-tuning yields \textbf{86.40\% Top-1 accuracy}---the
primary cross-domain performance indicator for the system---demonstrating that the
four-graph co-inference architecture captures generalisable molecular features beyond
the DOM training distribution.
An architectural extension, M11 (FCL-17 + cross-element interactions), achieves
91.50\% Top-1 on DOM (MRR\,=\,0.9427) and exposes a cross-class interference mechanism
detailed in Section~\ref{sec:run11_comparison}.
The key architectural drivers are the Spectral Mass-Difference Graph (contributing $-5.39$\,pp
when ablated) and FiLM instrument-type conditioning ($-4.20$\,pp), together encoding
inter-feature spectral context that is inaccessible to feature-independent scoring methods.

Element-class stratification
reveals near-perfect performance on CHO (99.8\%) and CHON (94.2\%) compounds, while
CHONS (18.9\%) and CHOPN (2.2\%) expose a deeper bottleneck: for phosphorus- and nitrogen-rich
compounds, the unrestricted candidate enumeration generates 100+ chemically equivalent
candidates, making discriminative ranking fundamentally impossible regardless of model
sophistication.  This finding reframes the challenge: for the dominant compound classes
in metabolomics (CHO, CHON), the \emph{ranking} problem is effectively solved;
the remaining frontier is \emph{candidate generation} for rare heteroatom-rich formulas.

Fine-tuning from M11 weights with a formula-bypass pathway achieves
\textbf{78.9\% Top-1 accuracy} on a held-out PFAS validation set
($n=375$; Top-3:\,90.1\%, Top-5:\,92.0\%, Top-10:\,96.3\%)---a 10.6\,pp
improvement over the prior partial-freeze baseline---demonstrating that the SMDG's
explicit CF$_2$ and HF transformation edges provide transferable PFAS-relevant
structure even from a model trained exclusively on DOM natural organic matter.

\begin{sloppypar}
The physics-informed design (differentiable Roepstorff isotopologue prediction,
Bethe--Weizs{\"a}cker mass-defect features, hard Seven Golden Rules constraints,
and physically consistent adduct-consistency penalties) demonstrates that domain
knowledge embedded in the neural architecture is synergistic with data-driven
learning rather than competing with it.
\end{sloppypar}

Future work should address: (i) CHOPN/CHONS-targeted candidate reduction
via class-conditional priors and adduct-pair disambiguation; (ii) positive-mode
and Orbitrap training to extend the model's instrument coverage;
(iii) integration of MS2 data through the DAG-GNN extension for compound classes
where MS1 alone is insufficient; (iv) improved calibration for risk-stratified
annotation workflows; and (v) application to ion mobility-integrated workflows
where the CCS scorer provides an additional discrimination axis.

MOSAIC-GNN establishes that deep multi-relational graph learning over spectral
mass-difference, isotopologue, formula-lattice, and adduct-consistency graphs
can close---and for the dominant compound classes entirely eliminate---the accuracy
gap between MS1-only and MS1+MS2 formula assignment.  We anticipate that this
framework will become a productive foundation for high-throughput untargeted
metabolomics annotation in environments where MS2 coverage is incomplete.
To make these capabilities accessible to non-programming users, we provide both
a web-based platform (Next.js frontend with FastAPI backend, supporting
drag-and-drop mzML and CSV upload, real-time progress tracking, and interactive
result visualisations) and a standalone cross-platform desktop application built
on Tkinter---deployable without any cloud dependency.  Both interfaces are
distributed together with the model weights and training code at the repository
available to all users.

% ═══════════════════════════════════════════════════════
\clearpage
\section*{Abbreviations}

% Single-spacing prevents the double-spaced body from pushing this over two pages
{\setstretch{1.0}\small
\begin{tabular}{l p{5.0cm} @{\hspace{1.5em}} l p{5.0cm}}
ADC  & Adduct Consistency Graph       & MRR  & Mean Reciprocal Rank \\
CCS  & Collision Cross-Section         & MS   & Mass Spectrometry \\
DBE  & Double Bond Equivalence         & MS1  & First-stage MS (precursor scan) \\
DOM  & Dissolved Organic Matter        & MS2  & Second-stage MS (fragmentation) \\
ECE  & Expected Calibration Error      & NIG  & Normal-Inverse-Gamma \\
ESI  & Electrospray Ionisation         & PFAS & Per/Polyfluoroalkyl Substances \\
ETD  & Element-Transformation Dict.    & pp   & Percentage Points \\
ETD  & Element-Transformation Dict.    & R-GCN & Relational Graph Conv.~Network \\
FCL  & Formula Candidate Lattice       & SMDG & Spectral Mass-Difference Graph \\
FiLM & Feature-wise Linear Modulation  & & \\
FT-ICR MS & FT Ion Cyclotron~Resonance MS & & \\
GAT  & Graph Attention Network         & & \\
GIN  & Graph Isomorphism Network       & & \\
GNPS & Global Natural Products Social  & & \\
GRL  & Gradient Reversal Layer         & & \\
HRMS & High-Resolution Mass Spectrometry & & \\
IHSS & Int.~Humic Substances Society   & & \\
IMS  & Ion Mobility Spectrometry       & & \\
ISO  & Isotopologue Pattern Graph      & & \\
KMD  & Kendrick Mass Defect            & & \\
\end{tabular}}

% ═══════════════════════════════════════════════════════
\section*{Data and Code Availability}

The MOSAIC-GNN model weights and inference code are available at
\url{https://github.com/Suppboat2544/MOSAIC-GNN_FormulaAssignment}.
The repository also includes a \textbf{web-based platform} (Next.js\,14 frontend +
FastAPI backend; drag-and-drop mzML/CSV upload; interactive per-class result histograms
and paginated formula tables) and a \textbf{standalone desktop application}
(cross-platform Tkinter GUI; no internet connection required), both designed for
non-programming users to run formula assignment directly on their instrument data.
Training datasets are not redistributed; the GNPS spectral library is publicly available
at \url{https://gnps.ucsd.edu}, and the DOM test standards (PPFA, SRFA2, SRFA3) are
available from the International Humic Substances Society
(\url{https://www.humicsubstances.org}).

\section*{Competing Interests}

The authors declare no competing interests.

\section*{Authors' Contributions}

S.K.\ conceived and designed the study, implemented the MOSAIC-GNN architecture and training
pipeline, performed all experiments and formal analyses, curated the data, created all
visualisations, and wrote the original manuscript draft.
N.S.-A.\ contributed to software engineering, system architecture design, web platform development, and manuscript review and editing.
S.Ch.\ contributed to data curation, formal analysis, and figure preparation.
N.D.\ contributed to cross-domain validation analysis, metabolomics domain expertise,
and manuscript review and editing.
G.Y.\ contributed to data curation, implementation support, and manuscript review.
M.F.\ directed the research, provided resources and supervision, acquired funding, and
contributed to manuscript review and editing.
All authors read and approved the final manuscript.

\section*{Acknowledgements}

This work was supported by JSPS KAKENHI Grant Numbers 23KK074 and 25K15476, and by
JST PRESTO (Grant Number JPMJPR24G7).
The authors thank the International Humic Substances Society (IHSS) for providing the PPFA,
SRFA2, and SRFA3 reference standards, and the GNPS team for maintaining the spectral library
used in cross-domain evaluation.

\section*{Authors' Information}

\textbf{Supaporn Klabklaydee}, Department of Civil and Environmental Engineering,
School of Environment and Society, Institute of Science Tokyo, Tokyo, Japan.

\textbf{Nopphakorn Subsa-Ard}, Department of Computer Engineering,
Faculty of Engineering, King Mongkut's University of Technology Thonburi,
Bangkok 10140, Thailand.

\textbf{Suppanut Chankhow}, Department of Civil and Environmental Engineering,
School of Environment and Society, Institute of Science Tokyo, Tokyo, Japan.

\textbf{Nattapak Darumas}, Department of Eye and Vision Sciences,
Institute of Life Course and Medical Sciences, University of Liverpool,
Liverpool, United Kingdom.

\textbf{Gu Yuheng}, Department of Civil and Environmental Engineering,
School of Environment and Society, Institute of Science Tokyo, Tokyo, Japan.

\textbf{Manabu Fujii}, Department of Civil and Environmental Engineering,
School of Environment and Society, Institute of Science Tokyo, Tokyo, Japan.

% ═══════════════════════════════════════════════════════
\begin{thebibliography}{99}

\bibitem{Kind2007}
Kind, T.; Fiehn, O.
Seven Golden Rules for heuristic filtering of molecular formulas obtained by accurate mass spectrometry.
\textit{BMC Bioinformatics} \textbf{2007}, \textit{8}, 105.
\doi{10.1186/1471-2105-8-105}

\bibitem{Hufsky2020}
Hufsky, F.; Böcker, S.
Mining molecular structure databases: Identification of small molecules based on fragmentation mass spectrometry data.
\textit{Mass Spectrometry Reviews} \textbf{2017}, \textit{36}(5), 624--633.

\bibitem{Duehrkop2019}
Dührkop, K.; Fleischauer, M.; Ludwig, M.; Aksenov, A. A.; Melnik, A. V.; Meusel, M.; Dorrestein, P. C.; Rousu, J.; Böcker, S.
SIRIUS 4: A rapid tool for turning tandem mass spectra into metabolite structure information.
\textit{Nature Methods} \textbf{2019}, \textit{16}(4), 299--302.
\doi{10.1038/s41592-019-0344-8}

\bibitem{Duehrkop2015}
Dührkop, K.; Shen, H.; Meusel, M.; Rousu, J.; Böcker, S.
Searching molecular structure databases with tandem mass spectra using CSI:FingerID.
\textit{Proceedings of the National Academy of Sciences} \textbf{2015}, \textit{112}(41), 12580--12585.

\bibitem{Dührkop2021}
Dührkop, K.; Nothias, L.-F.; Fleischauer, M.; et al.
Systematic classification of unknown metabolites using high-resolution fragmentation mass spectra.
\textit{Nature Biotechnology} \textbf{2021}, \textit{39}(4), 462--471.

\bibitem{Guo2023}
Guo, J.; et al.
BUDDY: Molecular formula discovery via bottom-up MS/MS interrogation.
\textit{Nature Methods} \textbf{2023}, \textit{20}(6), 881--890.
\doi{10.1038/s41592-023-01850-x}

\bibitem{Tolić2017}
Tolić, N.; et al.
Formularity: Software for Automated Formula Assignment of Natural and Other Organic Matter from Ultrahigh-Resolution Mass Spectra.
\textit{Analytical Chemistry} \textbf{2017}, \textit{89}(23), 12659--12665.

\bibitem{Wang2016}
Wang, M.; et al.
Sharing and community curation of mass spectrometry data with Global Natural Products Social Molecular Networking.
\textit{Nature Biotechnology} \textbf{2016}, \textit{34}(8), 828--837.
\doi{10.1038/nbt.3597}

\bibitem{Gilmer2017}
Gilmer, J.; Schütt, A. R.; Messmer, M. F.; Müller, K.-R.; Csányi, G.
Neural message passing for quantum chemistry.
\textit{Proceedings of Machine Learning Research} \textbf{2017}, \textit{70}, 1263--1272.

\bibitem{Xu2019}
Xu, K.; Hu, W.; Leskovec, J.; Jegelka, S.
How powerful are graph neural networks?
\textit{ICLR 2019}.

\bibitem{Reif2019}
Reif, M.; et al.
Graph neural networks for drug discovery.
\textit{Drug Discovery Today} \textbf{2019}.

\bibitem{Yang2019}
Yang, K.; et al.
Analyzing learned molecular representations for property prediction.
\textit{Journal of Chemical Information and Modeling} \textbf{2019}, \textit{59}(8), 3370--3388.

\bibitem{Bouwmeester2019}
Bouwmeester, R.; et al.
Comprehensive and empirical retention time prediction via machine learning for all ion monitoring metabolite identification.
\textit{Analytical Chemistry} \textbf{2019}, \textit{91}(5), 3694--3703.

\bibitem{Mesleh1996}
Mesleh, M. F.; Hunter, J. M.; Shvartsburg, A. A.; Schatz, G. C.; Jarrold, M. F.
Structural information from ion mobility measurements: Effects of the long-range potential.
\textit{Journal of Physical Chemistry} \textbf{1996}, \textit{100}(40), 16082--16086.

\bibitem{Perez2018}
Perez, E.; Strub, F.; de Vries, H.; Dumoulin, V.; Courville, A.
FiLM: Visual reasoning with a general conditioning layer.
\textit{Proceedings of the AAAI Conference on Artificial Intelligence} \textbf{2018}, \textit{32}(1).

\bibitem{Ganin2016}
Ganin, Y.; et al.
Domain-adversarial training of neural networks.
\textit{Journal of Machine Learning Research} \textbf{2016}, \textit{17}(59), 1--35.

\bibitem{Schütt2018}
Schütt, K. T.; Sauceda, H. E.; Kindermans, P.-J.; Tkatchenko, A.; Müller, K.-R.
SchNet – A deep learning architecture for molecules and materials.
\textit{Journal of Chemical Physics} \textbf{2018}, \textit{148}(24), 241722.

\bibitem{Aron2020}
Aron, A. T.; et al.
Reproducible molecular networking of untargeted mass spectrometry data using GNPS.
\textit{Nature Protocols} \textbf{2020}, \textit{15}(8), 2563--2584.

\bibitem{Stenson2003}
Stenson, A. C.; Landing, W. M.; Marshall, A. G.; Cooper, W. T.
Ionization and fragmentation of humic substances in electrospray ionization Fourier transform--ion cyclotron resonance mass spectrometry.
\textit{Analytical Chemistry} \textbf{2002}, \textit{74}(17), 4397--4409.

\bibitem{Roepstorff1984}
Roepstorff, P.; Fohlman, J.
Proposal for a common nomenclature for sequence ions in mass spectra of peptides.
\textit{Biomedical Mass Spectrometry} \textbf{1984}, \textit{11}(11), 601.

\end{thebibliography}

\section*{Supporting Information Available}
\noindent The Supporting Information is available free of charge on the ACS Publications website.
Mathematical specifications of all four graph representations
(Section~S1), full loss-function formulation with equations and weight
schedules (Section~S2), complete hyperparameter table (Section~S3),
visualization analysis of training dynamics and calibration
(Section~S4; Figures~S4.1--S4.6), extended GNPS cross-domain attribution
analysis with 12 figures (Section~S5), and H-series statistical
analysis with 15 inferential figures (Section~S6).

\end{document}
